\documentclass[12pt]{article}

\usepackage{sbc-template}
\usepackage{graphicx,url}
\usepackage[utf8]{inputenc}
\usepackage[brazil]{babel}
\usepackage{amsmath}
\usepackage{booktabs}
\usepackage{array}
\usepackage{enumitem}

\sloppy

\title{Impacto de Técnicas de Recuperação da Informação no Desempenho de Sistemas RAG para Tutores Inteligentes}

\author{Fabiana Marques Costa\inst{1}, Marcos Hideki Nishida Sakaguchi\inst{1}}

\address{Programa de Pós-Graduação em Ciência da Computação -- Instituto de Geociências e Ciências Exatas\\
Universidade Estadual Paulista ``Júlio de Mesquita Filho'' (UNESP)\\
Rio Claro -- SP -- Brasil
\email{fabiana.m.costa@unesp.br, marcos.sakaguchi@unesp.br}
}

\begin{document}

\maketitle

\begin{abstract}
This paper investigates the impact of Information Retrieval techniques on the performance of Retrieval-Augmented Generation systems for intelligent tutoring. Although RAG has been increasingly adopted to ground language models in external educational materials, the retrieval stage is often treated as an implementation detail rather than as a central experimental variable. This work proposes an evaluation protocol comparing sparse, dense, hybrid and reranked retrieval strategies in a RAG pipeline, using retrieval metrics and generation-oriented metrics. The study aims to clarify how retrieval quality affects answer correctness, faithfulness and pedagogical adequacy in educational question answering scenarios.
\end{abstract}

\begin{resumo}
Este artigo investiga o impacto de técnicas de Recuperação da Informação no desempenho de sistemas baseados em \textit{Retrieval-Augmented Generation} aplicados a tutores inteligentes. Embora arquiteturas RAG sejam utilizadas para ancorar modelos de linguagem em materiais educacionais externos, a etapa de recuperação é frequentemente tratada como detalhe de implementação. Este trabalho propõe um protocolo de avaliação comparando estratégias esparsas, densas, híbridas e com reranqueamento, utilizando métricas clássicas de RI e métricas orientadas à geração. O objetivo é analisar como a qualidade da recuperação influencia a correção, a fidelidade e a adequação pedagógica das respostas geradas.
\end{resumo}

\section{Introdução}

A incorporação de Inteligência Artificial (IA) em ambientes educacionais tem ampliado as possibilidades de personalização, acompanhamento e mediação do processo de aprendizagem. Sistemas Tutores Inteligentes (\textit{Intelligent Tutoring Systems} -- ITS) representam uma das linhas mais tradicionais desse campo, buscando adaptar conteúdos, feedbacks e estratégias pedagógicas ao perfil do estudante \cite{woolf2009}. Entretanto, a construção desses sistemas historicamente depende de modelos explícitos de domínio, regras pedagógicas e estruturas de conhecimento que podem ser custosas de manter, especialmente em contextos nos quais os conteúdos se modificam com frequência.

Mais recentemente, Modelos de Linguagem de Grande Escala (\textit{Large Language Models} -- LLMs) passaram a ser explorados como interfaces conversacionais para apoio ao ensino, devido à sua capacidade de gerar explicações em linguagem natural, adaptar exemplos e responder perguntas abertas. Apesar desse potencial, o uso isolado de LLMs em cenários educacionais apresenta riscos relevantes. Entre eles estão a geração de respostas não fundamentadas, a apresentação de informações incorretas com aparência de confiança e a dificuldade de garantir alinhamento ao material instrucional adotado por uma disciplina ou instituição. Em aplicações de tutoria, tais problemas são particularmente críticos, pois respostas incorretas podem induzir concepções equivocadas e comprometer a aprendizagem.

Nesse cenário, a arquitetura de \textit{Retrieval-Augmented Generation} (RAG) surge como alternativa promissora ao combinar recuperação de conhecimento externo com geração de linguagem natural \cite{lewis2020}. Em um pipeline RAG, documentos educacionais são indexados, consultas dos estudantes são usadas para recuperar trechos relevantes e o modelo gerador recebe esses trechos como contexto para produzir a resposta. Assim, o sistema deixa de depender exclusivamente do conhecimento paramétrico do LLM e passa a operar sobre uma base documental controlada, atualizável e potencialmente verificável.

Contudo, a efetividade de um sistema RAG depende diretamente da qualidade do componente de Recuperação da Informação (RI). Caso o recuperador selecione trechos irrelevantes, incompletos ou semanticamente distantes da dúvida do estudante, o gerador poderá produzir respostas vagas, incorretas ou pouco alinhadas ao material de referência. Por outro lado, uma recuperação de alta qualidade tende a fornecer evidências mais adequadas, favorecendo respostas mais fiéis ao corpus, mais úteis pedagogicamente e menos suscetíveis a alucinações.

Apesar dessa dependência, muitos trabalhos em RAG enfatizam o modelo generativo, o prompt ou a etapa de pós-processamento, enquanto a escolha da técnica de recuperação é frequentemente apresentada de forma secundária. Essa lacuna é ainda mais relevante em contextos educacionais, nos quais a avaliação não deve considerar apenas precisão factual, mas também clareza, coerência didática, adequação ao nível do estudante, rastreabilidade da fonte e capacidade de apoiar a aprendizagem. Dessa forma, torna-se necessário investigar a RI como variável central no desempenho de sistemas RAG voltados à tutoria inteligente.

Este artigo propõe um protocolo experimental para analisar o impacto de diferentes técnicas de RI em sistemas RAG educacionais. O foco recai sobre a comparação entre abordagens esparsas, densas, híbridas e com reranqueamento, avaliadas em duas camadas complementares: (i) a qualidade da recuperação, por meio de métricas clássicas de RI; e (ii) a qualidade da resposta gerada, por meio de métricas específicas para RAG e critérios pedagógicos. As principais contribuições esperadas deste trabalho são: 

\begin{enumerate}[label=(\roman*)]
    \item sistematizar o papel da Recuperação da Informação como componente determinante em tutores inteligentes baseados em RAG;
    \item propor um protocolo de avaliação em duas camadas, distinguindo métricas de recuperação e métricas de geração;
    \item comparar estratégias de recuperação esparsa, densa, híbrida e reranqueada em cenários de perguntas educacionais;
    \item analisar a relação entre métricas clássicas de RI e métricas de qualidade final das respostas geradas.
\end{enumerate}

\section{Fundamentação Teórica e Trabalhos Relacionados}

\subsection{Sistemas Tutores Inteligentes e IA na Educação}

Sistemas Tutores Inteligentes buscam reproduzir, em algum grau, características de um processo de tutoria individualizada. Em geral, esses sistemas são compostos por um modelo do domínio, responsável por representar o conhecimento a ser ensinado; um modelo do estudante, que descreve o estado de aprendizagem, dificuldades e progresso; um modelo pedagógico, que define estratégias de intervenção; e uma interface de interação \cite{woolf2009}. Essa estrutura permite que o sistema ofereça feedbacks personalizados, selecione atividades adequadas e acompanhe a evolução do estudante. Revisões recentes sobre chatbots educacionais também indicam a ampliação do uso de agentes conversacionais como suporte ao aprendizado, ainda que persistam desafios de avaliação, confiabilidade e integração pedagógica \cite{kuhail2023}.

A literatura em IA na Educação tem mostrado que sistemas educacionais inteligentes podem contribuir para aprendizagem ativa, prática guiada e feedback imediato. Entretanto, a construção de tutores tradicionais exige modelagem detalhada do domínio e engenharia de conhecimento, o que dificulta sua escalabilidade para múltiplas disciplinas, cursos e materiais. Além disso, quando o conteúdo muda, o sistema precisa ser atualizado manualmente, o que limita sua aplicabilidade em ambientes de ensino dinâmicos.

Com a popularização dos LLMs, novas possibilidades de tutoria conversacional passaram a ser exploradas. Diferentemente de tutores baseados apenas em regras, LLMs conseguem produzir explicações variadas, reformular conceitos, responder dúvidas abertas e adaptar o tom da resposta. No entanto, a fluidez linguística não garante correção, completude ou alinhamento pedagógico. Para que um tutor baseado em LLM seja confiável, é necessário que suas respostas estejam ancoradas em fontes apropriadas e que o sistema seja capaz de recuperar conteúdos relevantes antes de gerar a explicação.

\subsection{Recuperação da Informação em Sistemas Baseados em Conhecimento}

A Recuperação da Informação é a área responsável por representar, indexar, buscar e ordenar documentos de acordo com sua relevância para uma consulta \cite{manning2008}. Em sistemas educacionais, a RI pode ser utilizada para localizar trechos de livros, slides, apostilas, normas, exercícios resolvidos ou bases de conhecimento que respondam a uma pergunta do estudante. A forma como esses trechos são selecionados influencia diretamente a qualidade da mediação realizada pelo sistema.

Abordagens clássicas de RI, como TF-IDF e BM25, operam principalmente com base em correspondência lexical. O BM25, em particular, permanece amplamente utilizado por sua eficiência, simplicidade e bom desempenho em cenários nos quais os termos da consulta aparecem de maneira explícita nos documentos \cite{robertson2009}. Essa característica é útil em bases educacionais que possuem terminologia técnica específica, nomes de conceitos, fórmulas ou definições formais.

Entretanto, métodos lexicais apresentam limitações quando a consulta e o documento expressam a mesma ideia com vocabulários distintos. Por exemplo, uma pergunta como ``por que a busca vetorial encontra textos parecidos mesmo sem palavras iguais?'' pode demandar a recuperação de trechos que expliquem similaridade semântica, embeddings ou espaços vetoriais, ainda que esses termos não apareçam na consulta do estudante. Para lidar com esse tipo de cenário, técnicas neurais de recuperação densa representam consultas e documentos em espaços vetoriais contínuos, permitindo comparar proximidade semântica \cite{karpukhin2020}.

A recuperação densa, contudo, também apresenta desafios. Seu desempenho depende da qualidade dos embeddings, do domínio de treinamento, da estratégia de segmentação dos documentos e dos parâmetros de indexação. Além disso, modelos densos podem recuperar trechos semanticamente semelhantes, mas insuficientemente específicos para responder uma pergunta que exige termos exatos, datas, nomes ou definições formais. Por esse motivo, abordagens híbridas têm sido exploradas para combinar a precisão lexical de métodos esparsos com a flexibilidade semântica de métodos densos.

\subsection{Retrieval-Augmented Generation}

A arquitetura RAG combina um mecanismo de recuperação com um modelo gerativo. Originalmente proposta por Lewis et al. \cite{lewis2020}, a abordagem utiliza uma memória não paramétrica externa para fornecer evidências ao modelo de linguagem durante a geração. Em vez de produzir respostas apenas a partir dos parâmetros aprendidos no treinamento, o modelo passa a condicionar sua saída a documentos recuperados em tempo de inferência.

Um pipeline RAG típico envolve cinco etapas: (i) coleta e preparação do corpus; (ii) segmentação dos documentos em trechos menores, ou \textit{chunks}; (iii) indexação dos trechos por um ou mais métodos de recuperação; (iv) recuperação dos \textit{top-k} trechos mais relevantes para uma consulta; e (v) geração da resposta com base no contexto recuperado. Embora a etapa final seja frequentemente a mais visível para o usuário, a etapa de recuperação define o espaço informacional a partir do qual a resposta será construída.

Em tutores inteligentes, essa arquitetura pode ser utilizada para ancorar respostas em materiais oficiais da disciplina. Por exemplo, uma base RAG pode conter capítulos de livros, notas de aula, listas de exercícios e rubricas de avaliação. Quando o estudante pergunta sobre um conceito, o sistema recupera trechos desses materiais e gera uma explicação contextualizada. Essa dinâmica favorece rastreabilidade, atualização e controle sobre as fontes, desde que o componente de recuperação seja suficientemente eficaz.

\subsection{Avaliação de Sistemas RAG}

Avaliar sistemas RAG é uma tarefa complexa porque o desempenho final depende da interação entre múltiplos componentes. Uma resposta ruim pode decorrer de recuperação inadequada, de contexto excessivamente longo, de falha do gerador, de prompt mal formulado ou de divergência entre a pergunta e o corpus. Por isso, avaliar apenas a resposta final pode esconder problemas específicos do recuperador.

Ferramentas como RAGAS propõem métricas para avaliação automatizada de sistemas RAG, considerando dimensões como relevância do contexto, fidelidade da resposta ao contexto e relevância da resposta para a pergunta \cite{ragas2023}. Trabalhos posteriores também destacam a necessidade de avaliação mais centrada no recuperador. Além disso, levantamentos recentes sobre avaliação de RAG destacam que a análise deve considerar tanto recuperação quanto geração \cite{yu2024survey}. Ferramentas como ARES também reforçam a avaliação por dimensões, incluindo relevância de contexto, fidelidade e relevância da resposta \cite{ares2023}. Salemi e Zamani \cite{salemi2024} argumentam que métricas tradicionais baseadas apenas em relevância de documentos podem ter correlação limitada com o desempenho final do RAG, sugerindo métodos que aproximem a avaliação da recuperação da utilidade dos documentos para a geração.

Essas discussões reforçam que a avaliação de RAG deve ser realizada em camadas. A primeira camada mede se o sistema recupera os documentos ou trechos esperados; a segunda verifica se a resposta gerada é correta, fiel, útil e pedagogicamente adequada. Em contextos educacionais, essa separação é essencial, pois um trecho pode ser relevante do ponto de vista informacional, mas ainda assim insuficiente para gerar uma explicação clara para determinado nível de estudante.

\subsection{Lacunas na Literatura}

A análise dos trabalhos relacionados indica três lacunas principais. A primeira é a ausência de comparações sistemáticas entre técnicas de recuperação em cenários educacionais baseados em RAG. Embora existam estudos sobre RAG em tarefas gerais de QA e estudos sobre tutores baseados em LLMs, ainda são menos frequentes análises que isolem a técnica de RI como variável experimental principal.

A segunda lacuna está relacionada à integração entre métricas clássicas de RI e métricas de qualidade gerativa. Métricas como Recall@k, MRR e nDCG permitem medir qualidade de ranqueamento, mas não necessariamente indicam se a resposta final será fiel, clara ou útil. Por outro lado, métricas de RAG capturam dimensões da geração, mas podem não explicar qual etapa do pipeline causou o erro.

A terceira lacuna refere-se ao domínio educacional. A avaliação de respostas em tutores inteligentes precisa considerar aspectos pedagógicos adicionais, como adequação ao nível do estudante, presença de explicação gradual, ausência de resposta direta quando o objetivo é estimular raciocínio e alinhamento ao material da disciplina. Dessa forma, a análise de RI em RAG educacional deve ir além de benchmarks genéricos e considerar também corpora e perguntas representativas de situações reais de aprendizagem.

\section{Formulação do Problema}

O problema investigado neste trabalho pode ser formulado da seguinte forma: dado um corpus educacional $C$, um conjunto de perguntas de estudantes $Q$ e um pipeline RAG $P$, em que medida diferentes técnicas de Recuperação da Informação alteram a qualidade das respostas geradas por um tutor inteligente?

Neste estudo, a técnica de recuperação é tratada como variável independente. O modelo gerador, o prompt, o corpus, o tamanho dos \textit{chunks}, o número de documentos recuperados e os parâmetros de geração devem ser mantidos constantes sempre que possível. Dessa forma, busca-se reduzir interferências externas e observar com maior clareza o efeito da estratégia de recuperação sobre o desempenho final.

A investigação é orientada pelas seguintes questões de pesquisa:

\begin{description}[leftmargin=0.5cm]
    \item[RQ1:] Como técnicas esparsas, densas, híbridas e reranqueadas diferem em métricas clássicas de Recuperação da Informação em corpora educacionais?
    \item[RQ2:] Melhorias nas métricas de recuperação estão associadas a melhorias na qualidade das respostas geradas pelo sistema RAG?
    \item[RQ3:] Quais tipos de pergunta educacional se beneficiam mais de recuperação lexical, semântica ou híbrida?
    \item[RQ4:] Em que situações uma recuperação com melhor ranqueamento não resulta necessariamente em melhor resposta pedagógica?
\end{description}

A partir dessas questões, formulam-se as seguintes hipóteses:

\begin{description}[leftmargin=0.5cm]
    \item[H1:] Estratégias híbridas tendem a apresentar melhor equilíbrio entre precisão lexical e cobertura semântica do que estratégias exclusivamente esparsas ou densas.
    \item[H2:] O uso de reranqueamento tende a melhorar a qualidade dos contextos fornecidos ao gerador, especialmente em perguntas que exigem seleção fina entre trechos semanticamente próximos.
    \item[H3:] Métricas de recuperação como Recall@k e nDCG@k apresentam correlação positiva, mas não perfeita, com métricas de qualidade da resposta final.
    \item[H4:] Perguntas conceituais abertas tendem a se beneficiar mais de recuperação densa ou híbrida, enquanto perguntas definicionais e terminológicas podem favorecer métodos esparsos.
\end{description}

\section{Metodologia Proposta}

\subsection{Visão Geral do Pipeline Experimental}

A metodologia proposta organiza o sistema em um pipeline modular de RAG. O corpus educacional é inicialmente coletado, normalizado e segmentado em trechos. Em seguida, cada trecho é indexado por diferentes estratégias de recuperação. Para cada pergunta do conjunto de avaliação, o sistema recupera os \textit{top-k} trechos mais relevantes e os encaminha ao mesmo modelo gerador, utilizando um prompt fixo. A resposta produzida é então avaliada por métricas automatizadas e, quando possível, por critérios humanos ou semiautomáticos de qualidade pedagógica.

A adoção de um pipeline modular é importante porque permite substituir apenas o componente de recuperação, mantendo os demais elementos constantes. Assim, é possível comparar o impacto das técnicas de RI sem confundir seus efeitos com mudanças no modelo gerador ou no desenho do prompt.

\subsection{Implementação Inicial da Pipeline de Ingestão e Recuperação}

Como primeira etapa experimental, foi implementada uma pipeline de ingestão e avaliação de recuperação separada da etapa generativa. Essa decisão metodológica permite isolar o comportamento do recuperador antes de introduzir variáveis associadas ao LLM, ao prompt e à avaliação da resposta final. A pipeline inicial organiza o processo em quatro fases: (i) preparação do conjunto de dados; (ii) ingestão dos textos no índice vetorial e esparso; (iii) execução das consultas para diferentes métodos de recuperação; e (iv) cálculo agregado das métricas de RI.

No experimento preliminar, utilizou-se o SciQ como benchmark controlado. Cada parágrafo de suporte (\textit{support}) foi tratado como documento da base de conhecimento, enquanto cada pergunta (\textit{question}) foi usada como consulta. A relação entre a pergunta e seu parágrafo de suporte foi armazenada como julgamento de relevância (\textit{qrel}). Essa separação evita que perguntas, respostas corretas ou distratores sejam inseridos no índice, reduzindo o risco de vazamento experimental. Na configuração utilizada, após a remoção de exemplos sem suporte textual suficiente, foram gerados 12.135 documentos únicos, 12.146 consultas e 12.146 julgamentos de relevância, dos quais 878 consultas pertencem ao conjunto de teste.

A camada de indexação foi construída sobre uma coleção dedicada no Qdrant, com dois tipos de representação para o mesmo documento: vetores densos, produzidos por um modelo de embeddings, e vetores esparsos BM25, produzidos pelo FastEmbed. A busca lexical, a busca densa e a busca híbrida são executadas sobre a mesma coleção, o que reduz diferenças causadas por variações de corpus ou pré-processamento. A busca híbrida utiliza fusão por RRF, conforme descrito anteriormente, combinando os rankings esparso e denso.

Essa etapa corresponde à primeira parte do estudo: avaliar se a infraestrutura de ingestão e recuperação é capaz de localizar o trecho correto antes da geração. A segunda parte será composta pela avaliação RAG completa. Nessa fase, os trechos recuperados serão enviados a um modelo gerador com prompt fixo, e as respostas serão avaliadas com métricas do RAGAS, incluindo fidelidade ao contexto, relevância da resposta e relevância do contexto. Assim, a análise final poderá distinguir falhas da recuperação, falhas da geração e casos em que a recuperação é suficiente, mas o gerador ainda produz uma resposta inadequada.

\subsection{Preparação do Corpus}

O corpus deve ser composto por materiais educacionais representativos do domínio analisado. Em um cenário de aplicação real, esse corpus pode incluir apostilas, livros-texto, notas de aula, listas de exercícios, transcrições de aulas, documentos institucionais e respostas comentadas. Para fins de comparação controlada, também podem ser utilizados datasets públicos de question answering e recuperação, como SQuAD, HotpotQA, SciQ, ARC e subconjuntos do BEIR \cite{squad2016,hotpotqa2018,sciq2017,arc2018,thakur2021beir}.

A preparação envolve conversão dos documentos para texto, remoção de ruídos, preservação de metadados e segmentação em \textit{chunks}. A segmentação é uma etapa crítica, pois trechos muito longos podem introduzir ruído no contexto, enquanto trechos muito curtos podem perder informação necessária para responder à pergunta. Para cada \textit{chunk}, recomenda-se armazenar metadados como documento de origem, seção, página, tópico, data e nível educacional estimado.

\subsection{Estratégias de Recuperação Avaliadas}

Quatro estratégias principais são propostas para comparação:

\begin{enumerate}[label=(\alph*)]
    \item \textbf{BM25}: recuperação esparsa baseada em correspondência lexical, usada como baseline clássico;
    \item \textbf{Recuperação densa}: busca vetorial baseada em embeddings de consultas e documentos;
    \item \textbf{Recuperação híbrida}: combinação entre escores esparsos e densos, por fusão de rankings ou normalização de escores;
    \item \textbf{Recuperação híbrida com reranqueamento}: aplicação de um modelo reranqueador sobre os candidatos recuperados inicialmente.
\end{enumerate}

A estratégia híbrida pode ser implementada por \textit{Reciprocal Rank Fusion} (RRF), que combina listas ranqueadas atribuindo maior peso a documentos bem posicionados em diferentes recuperadores. A pontuação RRF de um documento $d$ pode ser expressa como:

\begin{equation}
RRF(d) = \sum_{r \in R} \frac{1}{k + rank_r(d)}
\end{equation}

em que $R$ representa o conjunto de rankings, $rank_r(d)$ é a posição do documento $d$ no ranking $r$ e $k$ é uma constante de suavização. Essa abordagem é útil por não exigir que os escores de BM25 e embeddings estejam na mesma escala.

\subsection{Arquitetura RAG}

A arquitetura RAG utilizada na avaliação deve manter constante o modelo gerador e o prompt. O prompt deve instruir o modelo a responder apenas com base nos trechos recuperados, indicar quando o contexto for insuficiente e evitar inventar informações não presentes no material. Essa decisão é especialmente importante em tutores inteligentes, pois respostas especulativas podem prejudicar o processo de aprendizagem.

Uma formulação simplificada do prompt pode incluir: (i) papel do sistema como tutor educacional; (ii) pergunta do estudante; (iii) contexto recuperado; (iv) instruções para explicar de forma didática; e (v) orientação para citar ou mencionar a base do conteúdo quando possível. O mesmo prompt deve ser utilizado em todas as condições experimentais para que a comparação reflita prioritariamente o efeito da recuperação.

\subsection{Avaliação da Recuperação}

A primeira camada de avaliação mede o desempenho do recuperador de forma independente da geração. Para isso, é necessário um conjunto de relevância, também conhecido como \textit{qrels}, indicando quais documentos ou trechos são relevantes para cada pergunta. Quando não houver anotação manual, pode-se construir um conjunto inicial a partir de perguntas sintéticas geradas com base nos próprios documentos, seguido de revisão humana em amostra.

As métricas propostas incluem Precision@k, Recall@k, Hit Rate@k, Mean Reciprocal Rank (MRR), Mean Average Precision (MAP) e Normalized Discounted Cumulative Gain (nDCG). A Precision@k estima a proporção de documentos relevantes entre os $k$ primeiros resultados:

\begin{equation}
Precision@k = \frac{|Rel_k|}{k}
\end{equation}

em que $Rel_k$ representa o conjunto de documentos relevantes recuperados até a posição $k$. O Recall@k mede a proporção de documentos relevantes existentes que foram recuperados:

\begin{equation}
Recall@k = \frac{|Rel_k|}{|Rel|}
\end{equation}

A MRR considera a posição do primeiro documento relevante, sendo adequada quando uma única evidência principal é suficiente para responder à pergunta. Já a nDCG considera a posição e o grau de relevância dos documentos, sendo útil quando há múltiplos trechos parcialmente relevantes.

\subsection{Avaliação da Resposta Gerada}

A segunda camada de avaliação mede a qualidade da resposta produzida pelo RAG. Para isso, propõe-se o uso de métricas automatizadas específicas para RAG, como as disponíveis no RAGAS \cite{ragas2023}. Entre as dimensões avaliadas estão: relevância do contexto, fidelidade da resposta ao contexto, relevância da resposta para a pergunta e similaridade com uma resposta de referência quando disponível.

Além das métricas automatizadas, recomenda-se uma rubrica pedagógica complementar. Essa rubrica pode avaliar: correção conceitual, clareza da explicação, adequação ao nível do estudante, uso apropriado de exemplos, ausência de informações não fundamentadas e capacidade de orientar o estudante sem simplesmente fornecer uma resposta final em atividades que exigem raciocínio. A Tabela~\ref{tab:metricas} sintetiza as duas camadas de avaliação.

\begin{table}[ht]
\centering
\caption{Camadas de avaliação propostas para o pipeline RAG educacional.}
\label{tab:metricas}
\begin{tabular}{p{3.2cm} p{5.0cm} p{5.0cm}}
\toprule
\textbf{Camada} & \textbf{Objetivo} & \textbf{Métricas ou critérios} \\
\midrule
Recuperação & Medir se os trechos corretos foram recuperados e bem ranqueados & Precision@k, Recall@k, Hit Rate@k, MRR, MAP, nDCG \\
Geração & Medir se a resposta é correta, fiel e útil ao estudante & Faithfulness, answer relevance, context relevance, correção conceitual, clareza e adequação pedagógica \\
\bottomrule
\end{tabular}
\end{table}

\subsection{Datasets e Cenários de Avaliação}

A avaliação pode combinar datasets públicos e um corpus educacional real. Datasets como SQuAD e HotpotQA permitem testar respostas extraídas de textos e perguntas que exigem múltiplas evidências \cite{squad2016,hotpotqa2018}. SciQ e ARC são úteis por conterem perguntas de ciências em formato educacional \cite{sciq2017,arc2018}. O BEIR, por sua vez, permite avaliar robustez em diferentes tarefas de recuperação \cite{thakur2021beir}. Estudos recentes sobre RAG educacional também reforçam a importância de considerar discrepâncias entre documentos instrucionais e conhecimento paramétrico do modelo, especialmente em ambientes escolares e universitários \cite{li2025educationrag,zheng2024edukdqa}. Quando o foco for educação brasileira, também podem ser construídos subconjuntos a partir de questões públicas do ENEM, desde que respeitadas as condições de uso e os critérios de curadoria.

A combinação de benchmarks públicos e corpus próprio é importante porque os benchmarks oferecem comparabilidade, enquanto o corpus real aproxima o experimento de uma aplicação de tutoria inteligente. Em ambos os casos, recomenda-se classificar as perguntas por tipo, por exemplo: definicionais, conceituais, procedimentais, comparativas, aplicadas e de múltiplas evidências. Essa classificação permite analisar quais técnicas de recuperação são mais adequadas para cada tipo de demanda educacional.

\section{Protocolo Experimental}

O protocolo experimental deve seguir uma sequência reprodutível. Primeiro, o corpus é processado e segmentado. Em seguida, são construídos índices separados para cada estratégia de recuperação. Depois, cada pergunta é executada em todas as condições experimentais. Para cada execução, devem ser armazenados: pergunta, técnica de recuperação, documentos recuperados, posições no ranking, escores, prompt final, resposta gerada, tempo de execução e métricas calculadas.

A Tabela~\ref{tab:condicoes} apresenta as condições experimentais sugeridas.

\begin{table}[ht]
\centering
\caption{Condições experimentais propostas.}
\label{tab:condicoes}
\begin{tabular}{p{2.8cm} p{4.8cm} p{5.4cm}}
\toprule
\textbf{Condição} & \textbf{Recuperador} & \textbf{Objetivo da comparação} \\
\midrule
C1 & BM25 & Estabelecer baseline lexical eficiente e interpretável \\
C2 & Embeddings densos & Avaliar ganhos de similaridade semântica \\
C3 & BM25 + denso & Medir complementaridade entre recuperação lexical e semântica \\
C4 & Híbrido + reranker & Verificar impacto de refinamento no ranqueamento final \\
\bottomrule
\end{tabular}
\end{table}

Para reduzir variabilidade, os seguintes controles devem ser aplicados: mesmo conjunto de perguntas, mesmo corpus, mesmo tamanho de \textit{chunk}, mesmo valor de $k$, mesmo modelo gerador, mesma temperatura de geração e mesmo prompt. Quando houver limitações computacionais, recomenda-se utilizar uma amostra estratificada de perguntas por tipo, preservando a diversidade de dificuldades e temas.

A análise estatística pode ser conduzida em duas etapas. Na primeira, compara-se o desempenho médio das estratégias de recuperação por métrica. Na segunda, calcula-se a correlação entre métricas de RI e métricas de geração, por exemplo utilizando Spearman ou Kendall. Essa análise permite verificar se melhorias em Recall@k ou nDCG@k se traduzem em maior fidelidade, relevância e qualidade pedagógica das respostas.

\section{Resultados Iniciais e Plano de Análise}

\subsection{Resultados Iniciais da Camada de Recuperação}

Os primeiros resultados foram obtidos com a camada de ingestão e recuperação sobre o conjunto de teste do SciQ. Foram comparadas três condições: BM25, recuperação densa e recuperação híbrida. A Tabela~\ref{tab:resultados-iniciais-sciq} apresenta os resultados agregados para Hit Rate@1, Hit Rate@5, Hit Rate@10 e MRR@10.

\begin{table}[ht]
\centering
\caption{Resultados iniciais de recuperação no SciQ, usando 878 consultas de teste.}
\label{tab:resultados-iniciais-sciq}
\begin{tabular}{lcccc}
\toprule
\textbf{Método} & \textbf{Hit@1} & \textbf{Hit@5} & \textbf{Hit@10} & \textbf{MRR@10} \\
\midrule
BM25 & 0,674 & 0,905 & 0,952 & 0,776 \\
Denso & 0,468 & 0,764 & 0,843 & 0,594 \\
Híbrido & 0,639 & 0,923 & 0,954 & 0,761 \\
\bottomrule
\end{tabular}
\end{table}

Os resultados indicam que o BM25 constitui uma baseline forte para o SciQ, especialmente na primeira posição do ranking. Esse comportamento é coerente com a natureza do dataset, pois muitas perguntas científicas compartilham termos técnicos ou expressões próximas às encontradas no parágrafo de suporte. A recuperação densa apresentou desempenho inferior nas métricas avaliadas, sugerindo que, nesse primeiro cenário, a similaridade semântica genérica não superou a correspondência lexical direta.

A estratégia híbrida apresentou o melhor Hit@5 e Hit@10, ainda que tenha ficado abaixo do BM25 em Hit@1 e MRR@10. Esse resultado sugere que a fusão entre recuperação lexical e densa aumenta a probabilidade de conter o documento correto no conjunto de contexto, mas nem sempre posiciona esse documento na primeira colocação. Para um sistema RAG, esse comportamento é relevante: quando o modelo gerador recebe múltiplos trechos, um maior Recall@k pode favorecer a presença da evidência necessária, mas a posição e a quantidade de ruído no contexto ainda podem afetar a resposta final.

Não foram observados indícios fortes de vazamento na preparação inicial. Cada consulta possui um único julgamento de relevância, os documentos indexados correspondem aos parágrafos de suporte e apenas uma pequena quantidade de documentos relevantes do teste também aparece associada a outros splits. Ainda assim, a interpretação dos resultados deve considerar uma limitação importante: em alguns casos, documentos semanticamente equivalentes ao suporte esperado podem ser recuperados, mas contam como erro por não possuírem o mesmo identificador do \textit{qrel}. Dessa forma, as métricas podem subestimar parcialmente a utilidade real de certos resultados recuperados.

\subsection{Plano de Análise da Etapa RAG}

Espera-se que estratégias híbridas apresentem desempenho mais estável do que estratégias isoladas, especialmente em corpora educacionais heterogêneos. A recuperação lexical tende a ser competitiva em perguntas que utilizam termos técnicos presentes no material, enquanto a recuperação densa tende a ser vantajosa quando o estudante formula a dúvida com vocabulário informal ou diferente do texto original. O reranqueamento, por sua vez, deve contribuir em cenários nos quais os primeiros candidatos recuperados são semanticamente próximos, mas variam quanto à utilidade real para responder à pergunta.

Também é esperado que as métricas de RI apresentem correlação positiva com métricas de RAG, mas que essa correlação não seja perfeita. Um alto Recall@k pode indicar que o trecho relevante está presente entre os documentos recuperados, mas não garante que ele esteja em posição privilegiada ou que o modelo gerador o utilize corretamente. Da mesma forma, uma resposta pode ser linguisticamente bem formulada mesmo quando o contexto recuperado é insuficiente, o que reforça a importância de métricas de fidelidade e avaliação pedagógica.

A análise qualitativa dos erros deve complementar os resultados quantitativos. Entre os erros esperados estão: recuperação de trechos genericamente relacionados, mas incapazes de responder à pergunta; seleção de trechos corretos, porém incompletos; conflito entre conhecimento paramétrico do LLM e material recuperado; e respostas que fornecem a solução final sem promover raciocínio. Essa análise é relevante para compreender não apenas qual técnica funciona melhor, mas por que determinadas estratégias falham em situações educacionais específicas.

\section{Ameaças à Validade}

Algumas ameaças à validade devem ser consideradas. A primeira está relacionada à construção do corpus. Se os documentos forem pouco representativos, mal segmentados ou excessivamente ruidosos, os resultados podem refletir problemas de preparação dos dados e não necessariamente limitações das técnicas de recuperação. A segunda ameaça envolve a anotação de relevância. Em domínios educacionais, diferentes trechos podem ser parcialmente úteis, o que torna a definição de relevância mais subjetiva.

Outra ameaça refere-se ao modelo gerador. Mesmo com recuperação adequada, o LLM pode ignorar parte do contexto, sintetizar informações incorretamente ou produzir explicações pouco didáticas. Por isso, a avaliação deve separar falhas de recuperação e falhas de geração sempre que possível. Além disso, o uso de métricas automatizadas baseadas em LLMs avaliadores pode introduzir vieses, exigindo validação com amostras avaliadas por humanos.

A validade externa também deve ser observada. Resultados obtidos em datasets públicos de QA podem não se transferir diretamente para disciplinas específicas, como programação, matemática, química ou ciências agrárias. Portanto, recomenda-se que o protocolo seja aplicado tanto em benchmarks padronizados quanto em um corpus educacional real, de modo a equilibrar comparabilidade e relevância prática.

\section{Considerações Finais}

Este artigo discutiu o papel da Recuperação da Informação no desempenho de sistemas RAG voltados a tutores inteligentes e propôs um protocolo experimental para avaliar esse impacto. A principal premissa é que a recuperação não deve ser tratada apenas como uma etapa auxiliar do pipeline, mas como componente central que condiciona a qualidade, a fidelidade e a utilidade pedagógica das respostas geradas.

A proposta organiza a avaliação em duas camadas complementares: uma camada de RI, baseada em métricas como Precision@k, Recall@k, MRR, MAP e nDCG; e uma camada de geração, baseada em métricas específicas para RAG e critérios pedagógicos. Essa separação permite diagnosticar se falhas observadas na resposta final decorrem de problemas na recuperação, na geração ou na interação entre ambas.

Como etapa inicial, foi implementada e executada a camada de ingestão e recuperação sobre o SciQ, permitindo comparar BM25, recuperação densa e recuperação híbrida antes da introdução do modelo gerador. Os resultados preliminares indicam que o BM25 é uma baseline forte nesse benchmark, enquanto a estratégia híbrida aumenta levemente a cobertura nos maiores valores de $k$. Como próximos passos, pretende-se executar a segunda camada do pipeline, na qual os contextos recuperados serão utilizados por um modelo gerador e avaliados com RAGAS e critérios pedagógicos. Espera-se que os resultados contribuam para o desenvolvimento de tutores inteligentes mais confiáveis, interpretáveis e alinhados aos materiais instrucionais, reforçando a importância da Recuperação da Informação como eixo metodológico em sistemas educacionais baseados em RAG.


\begin{thebibliography}{99}

\bibitem[ARC 2018]{arc2018}
Clark, P., Cowhey, I., Etzioni, O., Khot, T., Sabharwal, A., Schoenick, C. and Tafjord, O. (2018). ``Think You Have Solved Question Answering? Try ARC, the AI2 Reasoning Challenge''. \textit{arXiv preprint arXiv:1803.05457}.

\bibitem[ARES 2023]{ares2023}
Saad-Falcon, J., Khattab, O., Potts, C. and Zaharia, M. (2023). ``ARES: An Automated Evaluation Framework for Retrieval-Augmented Generation Systems''. \textit{arXiv preprint arXiv:2311.09476}.

\bibitem[Es et al. 2023]{ragas2023}
Es, S., James, J., Espinosa-Anke, L. and Schockaert, S. (2023). ``RAGAS: Automated Evaluation of Retrieval Augmented Generation''. \textit{arXiv preprint arXiv:2309.15217}.

\bibitem[HotpotQA 2018]{hotpotqa2018}
Yang, Z., Qi, P., Zhang, S., Bengio, Y., Cohen, W. W., Salakhutdinov, R. and Manning, C. D. (2018). ``HotpotQA: A Dataset for Diverse, Explainable Multi-hop Question Answering''. In \textit{Proceedings of the 2018 Conference on Empirical Methods in Natural Language Processing}, p. 2369--2380.

\bibitem[Karpukhin et al. 2020]{karpukhin2020}
Karpukhin, V., O{\u{g}}uz, B., Min, S., Lewis, P., Wu, L., Edunov, S., Chen, D. and Yih, W. (2020). ``Dense Passage Retrieval for Open-Domain Question Answering''. In \textit{Proceedings of the 2020 Conference on Empirical Methods in Natural Language Processing}, p. 6769--6781.

\bibitem[Kuhail et al. 2023]{kuhail2023}
Kuhail, M. A., Alturki, N., Alramlawi, S. and Alhejori, K. (2023). ``Interacting with Educational Chatbots: A Systematic Review''. \textit{Education and Information Technologies}, v. 28, p. 973--1018.

\bibitem[Lewis et al. 2020]{lewis2020}
Lewis, P., Perez, E., Piktus, A., Petroni, F., Karpukhin, V., Goyal, N., K{\"u}ttler, H., Lewis, M., Yih, W., Rockt{\"a}schel, T., Riedel, S. and Kiela, D. (2020). ``Retrieval-Augmented Generation for Knowledge-Intensive NLP Tasks''. In \textit{Advances in Neural Information Processing Systems}, v. 33, p. 9459--9474.

\bibitem[Li et al. 2025]{li2025educationrag}
Li, Z. et al. (2025). ``Retrieval-Augmented Generation for Educational Application: A Systematic Survey''. \textit{Computers and Education: Artificial Intelligence}.

\bibitem[Manning et al. 2008]{manning2008}
Manning, C. D., Raghavan, P. and Sch{\"u}tze, H. (2008). \textit{Introduction to Information Retrieval}. Cambridge University Press.

\bibitem[Rajpurkar et al. 2016]{squad2016}
Rajpurkar, P., Zhang, J., Lopyrev, K. and Liang, P. (2016). ``SQuAD: 100,000+ Questions for Machine Comprehension of Text''. In \textit{Proceedings of the 2016 Conference on Empirical Methods in Natural Language Processing}, p. 2383--2392.

\bibitem[Robertson and Zaragoza 2009]{robertson2009}
Robertson, S. and Zaragoza, H. (2009). ``The Probabilistic Relevance Framework: BM25 and Beyond''. \textit{Foundations and Trends in Information Retrieval}, v. 3, n. 4, p. 333--389.

\bibitem[Salemi and Zamani 2024]{salemi2024}
Salemi, A. and Zamani, H. (2024). ``Evaluating Retrieval Quality in Retrieval-Augmented Generation''. \textit{arXiv preprint arXiv:2404.13781}.

\bibitem[SciQ 2017]{sciq2017}
Welbl, J., Liu, N. F. and Gardner, M. (2017). ``Crowdsourcing Multiple Choice Science Questions''. In \textit{Proceedings of the 3rd Workshop on Noisy User-generated Text}, p. 94--106.

\bibitem[Thakur et al. 2021]{thakur2021beir}
Thakur, N., Reimers, N., R{\"u}ckl{\'e}, A., Srivastava, A. and Gurevych, I. (2021). ``BEIR: A Heterogeneous Benchmark for Zero-shot Evaluation of Information Retrieval Models''. \textit{arXiv preprint arXiv:2104.08663}.

\bibitem[Woolf 2009]{woolf2009}
Woolf, B. P. (2009). \textit{Building Intelligent Interactive Tutors: Student-Centered Strategies for Revolutionizing E-Learning}. Morgan Kaufmann.

\bibitem[Yu et al. 2024]{yu2024survey}
Yu, H., Gan, A., Zhang, K., Tong, S., Liu, Q. and Liu, Z. (2024). ``Evaluation of Retrieval-Augmented Generation: A Survey''. \textit{arXiv preprint arXiv:2405.07437}.

\bibitem[Zheng et al. 2024]{zheng2024edukdqa}
Zheng, T., Li, W., Bai, J., Wang, W. and Song, Y. (2024). ``Assessing the Robustness of Retrieval-Augmented Generation Systems in K-12 Educational Question Answering with Knowledge Discrepancies''. \textit{arXiv preprint arXiv:2412.08985}.

\end{thebibliography}

\end{document}
